\chapter*{Conclusion générale et Perspectives}
\addcontentsline{toc}{chapter}{Conclusion générale et Perspectives}
\markboth{Conclusion générale}{Conclusion générale}
\label{sec:conclusion}

\section{Rappel du contexte}

Un modèle HAR figé après l'entraînement initial tient mal en conditions
réelles : nouveaux porteurs, autres capteurs, activités non vues. Réapprendre
sans garde-fou contre l'oubli abîme ce qui marchait ; rester purement réactif
empêche d'alerter avant une transition. Ce mémoire s'est concentré sur ces
deux points : adaptation continue et anticipation légère.

\section{Ce que nous avons apporté}

L'encodeur \emph{IMUTransformerEncoder} (environ 531\,K paramètres, 4 blocs)
atteint 91,3~\% d'accuracy sur HAPT avec augmentation, sans pré-entraînement
massif.

Le tampon UWR priorise les fenêtres à forte entropie de prédiction. Avec
2\,000 exemples mémorisés, l'oubli passe d'environ 0,24 (finetuning naïf) à
0,09 sur le protocole à 10 sujets ; UWR dépasse EWC et iCaRL dans nos tests.

La tête d'anticipation LSTM prédit l'activité suivante à partir de $S=5$
fenêtres partielles. Le macro-F1 de validation se situe autour de 0,60--0,64 :
ce n'est pas encore un système d'alerte clinique, mais ce n'est plus du
hasard non plus (baseline majorité $\approx$0,04).

Pour le cross-dataset, le code livre un \texttt{DomainAdapter} affine
few-shot. Les chiffres mesurés en zéro-shot sont d'environ 0,46 sur WISDM et
0,12 sur PAMAP2. CORAL apparaît dans la bibliographie, pas comme perte
exécutée dans le dépôt.

Nous avons aussi ajouté un SSL SimCLR-style, un pipeline foundation-style
léger, une calibration de température, un bandit de seuil et un fall-risk
binaire (macro-F1 $\approx$0,88 sur les transitions 9--12). Le scénario
« repas » reste hors portée faute de labels cuisine. Un détecteur d'impact
hybride historique donne une AUC-ROC de 0,762 sur un proxy HAPT ; les chutes
réelles demanderaient MobiAct.

\section{Bilan chiffré}

\begin{table}[H]
\centering
\small
\begin{tabularx}{0.92\textwidth}{|Y|c|c|c|}
\hline
\textbf{Tâche} & \textbf{Baseline} & \textbf{Notre système} & \textbf{Gain} \\
\hline
Classification HAR       & 81,0~\% & 91,3~\% & +10,3 pts \\
CL user-incrémental        & 62,4~\% & 79,2~\% & +16,8 pts \\
CL class-incrémental       & 62,4~\% & 79,8~\% & +17,4 pts \\
Anticipation (macro-F1)    & —       & $\approx$0,60--0,64 & vs majorité $\approx$0,04 \\
Cross-dataset WISDM (F1)   & 0,838 (intra) & $\approx$0,46 & domain shift \\
Fall-risk (macro-F1)       & —       & $\approx$0,88 & transitions 9--12 \\
Détection impact (AUC)     & —       & 0,762   & — \\
\hline
\end{tabularx}
\caption{Bilan quantitatif des contributions}
\label{tab:bilan_conclusion}
\end{table}

\section{Limites}

Nous restons à environ 12 points sous l'entraînement joint en CL user. La
mémoire est bornée et les données n'arrivent qu'en séquence. L'anticipation
multi-classe reste fragile pour une alerte médicale. Les « chutes » du
fall-risk sont des transitions posturales, pas des chutes cliniques. SSL,
foundation-style et RL sont volontairement légers. Sans corpus cuisine, le
scénario repas de la fiche ne peut pas être entraîné.

\section{Perspectives}

À court terme, intégrer MobiAct, mieux gérer le déséquilibre des transitions
et, si possible, un jeu ADL cuisine. À moyen terme, exporter vers ONNX/TFLite,
tester OPPORTUNITY et pousser le SSL avec plus de données. Plus loin, un
agent qui composerait le tampon, un vrai backbone foundation
\parencite{yuan2024} branché sur UWR, ou une fusion IMU avec d'autres
modalités.

Le dépôt est rangé pour qu'on puisse reprendre le travail ou viser un
workshop. Les figures se régénèrent et les checkpoints évitent de tout
réentraîner. Une note courte centrée sur UWR, avec un protocole OCL-HAR
strict, serait un bon prochain pas.

\section{Retour sur les objectifs}

Le backbone tient ses promesses sur HAPT. Le CL est solide en user-incrémental
mais encore fragile en class-incrémental (F1 final faible à cause du
déséquilibre). L'anticipation et le fall-risk fonctionnent en démo. Le
cross-dataset montre clairement le domain shift ; l'adaptateur few-shot est
là pour l'atténuer. Les modules SSL / calibration / bandit répondent aux
orientations de la fiche sans prétendre à l'état de l'art industriel.

\section{Ce qu'on retiendrait pour un déploiement}

Un pré-entraînement supervisé soigné suffit souvent. Un tampon de 2\,000
fenêtres avec UWR vaut mieux qu'un finetuning naïf quand les utilisateurs
arrivent les uns après les autres. Pour alerter, le fall-risk binaire avec
seuil calibré est plus crédible que l'anticipation multi-classe seule. En
cross-domain, quelques labels cibles aident beaucoup : le zéro-shot ne
suffit pas. Les scripts et annexes détaillent comment reproduire ces
constats.
